% ---- begin paper/tables/t5_dosis.tex
\begin{table}[htbp]\centering
\caption{The Gradient Under Continuous Seismic Doses}
\label{tab:dose}
\begin{threeparttable}
\begin{tabular}{@{}p{170pt}*{4}{>{\centering\arraybackslash}p{67pt}}@{}}
\toprule
 & \multicolumn{4}{c}{Seismic dose of the municipality of selection} \\
\cmidrule(lr){2-5}
 & Affected zone (0/1) & PGA (z) & MMI (z) & $-$log dist.\ to epicenter (z) \\
 & (1) & (2) & (3) & (4) \\
\midrule
\panel{5}{Panel A. PHQ-4 score: sum of four depression--anxiety items (z)}
Exposed at 12--23 & $-$0.045 & $-$0.035 & $-$0.028 & $-$0.029 \\
 & (0.078) & (0.033) & (0.032) & (0.029) \\
\addlinespace[3pt]
Exposed at 24--35 & $-$0.147\sym{**} & $-$0.068\sym{**} & $-$0.062\sym{***} & $-$0.062\sym{**} \\
 & (0.069) & (0.027) & (0.023) & (0.025) \\
\addlinespace[3pt]
Exposed at 36--59 & $-$0.044 & $-$0.037 & $-$0.039 & $-$0.025 \\
 & (0.060) & (0.027) & (0.024) & (0.023) \\
\addlinespace
\panel{5}{Panel B. Positive GAD-2 screen: two anxiety items, subscore of 3 or more}
Exposed at 12--23 & $-$0.012 & $-$0.015 & $-$0.013 & $-$0.022\sym{*} \\
 & (0.033) & (0.014) & (0.013) & (0.013) \\
\addlinespace[3pt]
Exposed at 24--35 & $-$0.073\sym{**} & $-$0.039\sym{***} & $-$0.029\sym{**} & $-$0.036\sym{***} \\
 & (0.032) & (0.015) & (0.013) & (0.013) \\
\addlinespace[3pt]
Exposed at 36--59 & $-$0.016 & $-$0.016 & $-$0.010 & $-$0.014 \\
 & (0.028) & (0.013) & (0.010) & (0.011) \\
\midrule
Municipality fixed effects & Yes & Yes & Yes & Yes \\
Birth-cohort fixed effects & Yes & Yes & Yes & Yes \\
Pre-earthquake controls & Yes & Yes & Yes & Yes \\
2012 controls & Yes & Yes & Yes & Yes \\
Adolescents & 9,996 & 9,996 & 9,996 & 9,996 \\
Selection municipalities (clusters) & 116 & 116 & 116 & 116 \\
\bottomrule
\end{tabular}
\begin{wpnotes}
\textit{Notes:} Estimates of equation~(1) under four measures of the seismic dose, all with the controls of column 3 of Table~\ref{tab:main}. Column 1 repeats the official binary affected zone. Column 2 uses the peak ground acceleration: the strongest shaking of the ground during the earthquake, expressed as a share of the acceleration of gravity, measured at the centroid of the municipality of selection (from the ShakeMap of the U.S. Geological Survey) and standardized; column 3 an intensity scale from the same source; column 4 minus the log of the distance from the centroid to the epicenter. The continuous doses vary across the 116 municipalities, within affected regions, which addresses inference concerns with a regional treatment. \bindef{} \clusternote{} \starnote
\end{wpnotes}
\end{threeparttable}
\end{table}
% ---- end paper/tables/t5_dosis.tex
